\section{Analysis: Why Online Self-Distillation Works}
\label{sec:analysis}

\paragraph{The flywheel.} The method's components are individually simple but
mutually reinforcing: weight updates improve the model; forward re-evaluation lets
the improved model recover problems it failed on arrival ($84$ of $1{,}463$
memory entries in the fused run entered late this way); recovered problems become
demonstrations and distillation targets; and fresher targets keep the KL objective
near on-policy. A frozen-weights method (ICL) receives none of this---re-answering
with an unchanged model recovers almost nothing---and pure RL receives the reward
but not the dense token-level signal.

\paragraph{Distillation as supervision, memory as scaffolding.} Wherever ICL
hurts (HotpotQA, misleading neighbors), self-distillation still helps: the hint
channel injects only the model's own \emph{verified} answer for the \emph{same}
problem, a signal that cannot mislead. Conversely, wherever ICL is strong
(DDXPlus), consolidation into weights preserves the gains while freeing the
context window.

\paragraph{Stability.} The distillation loss (Eq.~\ref{eq:kl}) is anchored to a
frozen reference teacher and its targets are always oracle-verified text. There
is no advantage estimate, no clipped ratio, no moving critic---the failure
surfaces that destroyed REINFORCE++ on the long fused stream simply do not exist.
Across ${\sim}12{,}000$ streamed problems in this paper, online SDFT never
regressed catastrophically.

\paragraph{Predict-then-train integrity.} Because every reported number is
produced by a model that had never seen that problem---scored before storage,
demonstrations strictly from the past, hints strictly self-generated---the
streaming metric doubles as a continual-generalization measure. We release a
causal-ordering audit that verifies these invariants over the raw run logs
(Appendix~\ref{app:leakage}).
